% SPDX-License-Identifier: CC-BY-SA-4.0
% Author: Matthieu Perrin

\begingroup

\begin{exercice}[Révisions -- Le lemme de l'étoile] \label{exo:lexing/expressiveness/revision}

  Pour chacun des langages suivants sur l’alphabet $\Sigma = \{a,b\}$,
  déterminer s’il est rationnel. Justifier votre réponse :

  \begin{question}

  \item $\{ a^i b^j \mid i \not\equiv j \pmod{3} \}$.
  
  \item $\{ a^n b^m a^n \mid n,m \in \mathbb{N} \}$.
  
  \item $\{ xx \mid x \in \{a,b\}^\star \}$.

  \item $\{ a^{2^n} \mid n \in \mathbb{N} \}$.

  \end{question}

\end{exercice}

\begin{correction}

  \begin{question}
  \item Supposons (par l'absurde), que $L_1$ est rationnel.\\
    Alors $L_1$ vérifie le lemme de l'étoile, donc il existe $N \in \mathbb{N}$ tel que pour tout mot $u$ tel que $|u|\ge N$,
    il existe $x, y, z \in \Sigma^\star$ tels que $u=xyz$, $y\neq \emptyset$, $|xy| \le N$ et pour tout $i\in \mathbb{N}$, $xy^i z\in L_1$. \\
    \textbf{Posons $u = a^Nb^N$. } Clairement, on a $|u| = 2N \ge N$ et $u\in L_1$.\\
    Comme $|xy| \le N$, $xy$ n'est composé que de $a$, donc on a $x = a^{|x|}$, $y = a^{|y|} \neq \varepsilon$  et $z=a^{N-|x|-|y|}b^N$. \\
    \textbf{Posons $i = 0$. } On a $x y^i z = a^{N-|y|}b^N \in L_1$ d'après le lemme de l'étoile.\\
    Absurde (car $n = N-|y| \not\ge N = p$) donc $L_1$ n'est pas rationnel. 
  \item $L_2 = \mathcal{S}((aa)^\star \cdot (bb)^\star \mid a(aa)^\star \cdot b(bb)^\star )$ est rationnel.
  \item Supposons (par l'absurde), que $L_3$ est rationnel.\\
    Soit $N$ donné par le lemme de l'étoile. \\
    Posons $u = a^{{(N+1)}^2}$. On a bien $u \in L_3$, et $|u| \ge N$. \\
    Soit $x\cdot y \cdot z$ la décomposition donnée par le lemme de l'étoile. \\
    Comme $y\neq \varepsilon$ et $|y| \le |xy|\le N < N+1$, on a:
    $$\begin{array}{rcccl}
      0 &<& |y| &<& N+1\\
      (N+1)^2 &<& (N+1)^2 + |y| & < & (N+1)^2 + N+1\\
      (N+1)^2 &<& (N+1)^2 + |y| & < &  (N+1)(N+2)\\
      (N+1)^2 &<& |x\cdot y^2 \cdot z| & < & (N+2)^2\\
    \end{array}$$

    $(N+1)^2$ et $(N+2)^2$ sont deux carrés parfaits consécutifs, donc $|x\cdot y^2 \cdot z|$ n'est pas un carré parfait.
    Donc $x\cdot y^2 \cdot z \notin L_3$. Absurde, donc $L_3$ n'est pas rationnel. 
  \end{question}

\end{correction}

\endgroup
\endinput
